% --- Archon physics formalization source begin ---
% archon:physics
% archon:covers PhyXMiniProblems/problem_phyx_mini_0832.lean
% archon:source-report reports/phyx_mini/problem_phyx_mini_0832.source.json

\chapter{Physics problem phyx\_mini\_0832}
\label{ch:PhyXMiniProblems_problem_phyx_mini_0832}

\paragraph{Problem source.}
\#\# Question

What is the magnitude of the strength of the electric field at the position indicated by the dot?

\#\# Answer choices

- A: A: 1.3 \textbackslash{}times 10\textasciicircum{}\{5\}N/C

- B: B: 5.3 \textbackslash{}times 10\textasciicircum{}\{5\}N/C

- C: C: 1.0 \textbackslash{}times 10\textasciicircum{}\{5\}N/C

- D: D: 1.1 \textbackslash{}times 10\textasciicircum{}\{5\}N/C

\#\# Figure caption (auxiliary; use the image as primary evidence)

The image features an arrangement of point charges at the corners of a rectangle. Here's a detailed description:

- The top left corner has a circle labeled with a negative charge of -5.0 nC.
- The top right corner has a circle labeled with a positive charge of +10 nC.
- The bottom left corner has a circle labeled with a positive charge of +5.0 nC.
- The bottom right corner is marked with a small black dot but no charge.

The rectangle is oriented vertically with dashed lines connecting the charges, representing distances:
- The horizontal distance between the charges at the top is labeled as 2.0 cm.
- The vertical distance between the top and bottom charges is labeled as 4.0 cm.
- A horizontal dashed line also connects the bottom left corner to the black dot at the bottom right.

This setup likely represents an electric field or force vector problem.

\#\# Dataset metadata

Electrostatics; Physical Model Grounding Reasoning, Multi-Formula Reasoning

\paragraph{Recorded answer/context.}
D: D: 1.1 \textbackslash{}times 10\textasciicircum{}\{5\}N/C

\paragraph{Figure/image path.}
/root/proposal\_for\_physic/science-mango/phyx\_mini\_run/phyx\_data/test\_image/832.png

\paragraph{Formalization target.}
create a compiling Lean file with sorry bodies at `PhyXMiniProblems/problem\_phyx\_mini\_0832.lean`. The Lean declarations must preserve the physical quantities, dimensions or dimensional roles, figure labels, governing-law hypotheses, and final relation expressed by this problem.
Use Mathlib/Physlib names found through LeanExplore where available. If a physics API is missing, introduce faithful local abstractions rather than scalar placeholder aliases.

\begin{theorem}[Physics formalization target]
\label{thm:physics:phyx_mini_0832:target}
\lean{PhyXMiniProblems.ProblemPhyXMini0832.electric_field_magnitude_at_dot_selects_D}
\uses{def:physics:phyx-mini-0832:phyxminiproblems-problemphyxmini0832-rectangularpointchargesetup, def:physics:phyx-mini-0832:phyxminiproblems-problemphyxmini0832-matchesprimaryfigure832, def:physics:phyx-mini-0832:phyxminiproblems-problemphyxmini0832-usesstandardcoulombconstant, def:physics:phyx-mini-0832:phyxminiproblems-problemphyxmini0832-hasphysicalrectangularchargeparameters, def:physics:phyx-mini-0832:phyxminiproblems-problemphyxmini0832-satisfiescoulombfieldandsuperpositionlaws, def:physics:phyx-mini-0832:phyxminiproblems-problemphyxmini0832-isclosestdisplayedfieldstrength}
The assigned autoformalize agent should translate this physics problem into Lean declarations in the covered file, with theorem and lemma proof bodies written as `by sorry`.
\end{theorem}
\begin{proof}
This is an autoformalization task, not a proof task. Produce statements that can later be proved by the physics prover without weakening the physical model.
\end{proof}

% --- Archon declaration topology begin ---
\section{Declaration topology}
\begin{definition}[Length Quantity]
  \label{def:physics:phyx-mini-0832:phyxminiproblems-problemphyxmini0832-lengthquantity}
  \lean{PhyXMiniProblems.ProblemPhyXMini0832.LengthQuantity}
  A nonnegative physical length
\end{definition}
\begin{proof}
  This is the declared model definition of Length Quantity.
\end{proof}
\begin{definition}[Signed Charge Quantity]
  \label{def:physics:phyx-mini-0832:phyxminiproblems-problemphyxmini0832-signedchargequantity}
  \lean{PhyXMiniProblems.ProblemPhyXMini0832.SignedChargeQuantity}
  A signed electric charge, rather than a bare scalar charge alias
\end{definition}
\begin{proof}
  This is the declared model definition of Signed Charge Quantity.
\end{proof}
\begin{definition}[electric Field Strength Dimension]
  \label{def:physics:phyx-mini-0832:phyxminiproblems-problemphyxmini0832-electricfieldstrengthdimension}
  \lean{PhyXMiniProblems.ProblemPhyXMini0832.electricFieldStrengthDimension}
  The dimension M L T⁻² C⁻¹ of electric-field strength
\end{definition}
\begin{proof}
  This is the declared model definition of electric Field Strength Dimension.
\end{proof}
\begin{definition}[Electric Field Strength Quantity]
  \label{def:physics:phyx-mini-0832:phyxminiproblems-problemphyxmini0832-electricfieldstrengthquantity}
  \lean{PhyXMiniProblems.ProblemPhyXMini0832.ElectricFieldStrengthQuantity}
  \uses{def:physics:phyx-mini-0832:phyxminiproblems-problemphyxmini0832-electricfieldstrengthdimension}
  A nonnegative physical magnitude of electric-field strength
\end{definition}
\begin{proof}
  This is the declared model definition of Electric Field Strength Quantity.
\end{proof}
\begin{definition}[length Readout]
  \label{def:physics:phyx-mini-0832:phyxminiproblems-problemphyxmini0832-lengthreadout}
  \lean{PhyXMiniProblems.ProblemPhyXMini0832.lengthReadout}
  \uses{def:physics:phyx-mini-0832:phyxminiproblems-problemphyxmini0832-lengthquantity}
  Read a physical length in a selected unit
\end{definition}
\begin{proof}
  This is the declared model definition of length Readout.
\end{proof}
\begin{definition}[signed Charge In Coulombs]
  \label{def:physics:phyx-mini-0832:phyxminiproblems-problemphyxmini0832-signedchargeincoulombs}
  \lean{PhyXMiniProblems.ProblemPhyXMini0832.signedChargeInCoulombs}
  \uses{def:physics:phyx-mini-0832:phyxminiproblems-problemphyxmini0832-signedchargequantity}
  SI readout of a signed charge, in coulombs
\end{definition}
\begin{proof}
  This is the declared model definition of signed Charge In Coulombs.
\end{proof}
\begin{definition}[electric Field Strength In Newtons Per Coulomb]
  \label{def:physics:phyx-mini-0832:phyxminiproblems-problemphyxmini0832-electricfieldstrengthinnewtonspercoulomb}
  \lean{PhyXMiniProblems.ProblemPhyXMini0832.electricFieldStrengthInNewtonsPerCoulomb}
  \uses{def:physics:phyx-mini-0832:phyxminiproblems-problemphyxmini0832-electricfieldstrengthquantity}
  SI readout of electric-field strength, in newtons per coulomb
\end{definition}
\begin{proof}
  This is the declared model definition of electric Field Strength In Newtons Per Coulomb.
\end{proof}
\begin{definition}[Source Charge]
  \label{def:physics:phyx-mini-0832:phyxminiproblems-problemphyxmini0832-sourcecharge}
  \lean{PhyXMiniProblems.ProblemPhyXMini0832.SourceCharge}
  The three charged corners visible in the image
\end{definition}
\begin{proof}
  This is the declared model definition of Source Charge.
\end{proof}
\begin{definition}[Figure Corner]
  \label{def:physics:phyx-mini-0832:phyxminiproblems-problemphyxmini0832-figurecorner}
  \lean{PhyXMiniProblems.ProblemPhyXMini0832.FigureCorner}
  The four corners of the dashed rectangle
\end{definition}
\begin{proof}
  This is the declared model definition of Figure Corner.
\end{proof}
\begin{definition}[source Corner]
  \label{def:physics:phyx-mini-0832:phyxminiproblems-problemphyxmini0832-sourcecorner}
  \lean{PhyXMiniProblems.ProblemPhyXMini0832.sourceCorner}
  \uses{def:physics:phyx-mini-0832:phyxminiproblems-problemphyxmini0832-sourcecharge, def:physics:phyx-mini-0832:phyxminiproblems-problemphyxmini0832-figurecorner}
  The figure corner occupied by each named source charge
\end{definition}
\begin{proof}
  This is the declared model definition of source Corner.
\end{proof}
\begin{definition}[Rectangular Point Charge Figure]
  \label{def:physics:phyx-mini-0832:phyxminiproblems-problemphyxmini0832-rectangularpointchargefigure}
  \lean{PhyXMiniProblems.ProblemPhyXMini0832.RectangularPointChargeFigure}
  \uses{def:physics:phyx-mini-0832:phyxminiproblems-problemphyxmini0832-lengthquantity, def:physics:phyx-mini-0832:phyxminiproblems-problemphyxmini0832-signedchargequantity, def:physics:phyx-mini-0832:phyxminiproblems-problemphyxmini0832-sourcecharge, def:physics:phyx-mini-0832:phyxminiproblems-problemphyxmini0832-figurecorner, def:physics:phyx-mini-0832:phyxminiproblems-problemphyxmini0832-sourcecorner}
  Typed content of the primary raster. Its length and charge labels are physical quantities; their numerical readouts are imposed separately by MatchesPrimaryFigure832
\end{definition}
\begin{proof}
  This is the declared model definition of Rectangular Point Charge Figure.
\end{proof}
\begin{definition}[Rectangular Point Charge Setup]
  \label{def:physics:phyx-mini-0832:phyxminiproblems-problemphyxmini0832-rectangularpointchargesetup}
  \lean{PhyXMiniProblems.ProblemPhyXMini0832.RectangularPointChargeSetup}
  \uses{def:physics:phyx-mini-0832:phyxminiproblems-problemphyxmini0832-lengthquantity, def:physics:phyx-mini-0832:phyxminiproblems-problemphyxmini0832-signedchargequantity, def:physics:phyx-mini-0832:phyxminiproblems-problemphyxmini0832-electricfieldstrengthquantity, def:physics:phyx-mini-0832:phyxminiproblems-problemphyxmini0832-sourcecharge, def:physics:phyx-mini-0832:phyxminiproblems-problemphyxmini0832-rectangularpointchargefigure}
  The electrostatic system and its independent observable. In particular, resultantFieldStrengthAtDot is not defined from answer D or from any displayed numerical choice
\end{definition}
\begin{proof}
  This is the declared model definition of Rectangular Point Charge Setup.
\end{proof}
\begin{definition}[x Coordinate In Meters]
  \label{def:physics:phyx-mini-0832:phyxminiproblems-problemphyxmini0832-xcoordinateinmeters}
  \lean{PhyXMiniProblems.ProblemPhyXMini0832.xCoordinateInMeters}
  Horizontal coordinate in the metre-calibrated chart used for the figure
\end{definition}
\begin{proof}
  This is the declared model definition of x Coordinate In Meters.
\end{proof}
\begin{definition}[y Coordinate In Meters]
  \label{def:physics:phyx-mini-0832:phyxminiproblems-problemphyxmini0832-ycoordinateinmeters}
  \lean{PhyXMiniProblems.ProblemPhyXMini0832.yCoordinateInMeters}
  Vertical coordinate in the metre-calibrated chart used for the figure
\end{definition}
\begin{proof}
  This is the declared model definition of y Coordinate In Meters.
\end{proof}
\begin{definition}[resultant Magnitude In Newtons Per Coulomb At]
  \label{def:physics:phyx-mini-0832:phyxminiproblems-problemphyxmini0832-resultantmagnitudeinnewtonspercoulombat}
  \lean{PhyXMiniProblems.ProblemPhyXMini0832.resultantMagnitudeInNewtonsPerCoulombAt}
  \uses{def:physics:phyx-mini-0832:phyxminiproblems-problemphyxmini0832-rectangularpointchargesetup}
  Magnitude of the resultant vector-field readout at a spacetime point
\end{definition}
\begin{proof}
  This is the declared model definition of resultant Magnitude In Newtons Per Coulomb At.
\end{proof}
\begin{definition}[coulomb Field Vector At]
  \label{def:physics:phyx-mini-0832:phyxminiproblems-problemphyxmini0832-coulombfieldvectorat}
  \lean{PhyXMiniProblems.ProblemPhyXMini0832.coulombFieldVectorAt}
  \uses{def:physics:phyx-mini-0832:phyxminiproblems-problemphyxmini0832-signedchargeincoulombs, def:physics:phyx-mini-0832:phyxminiproblems-problemphyxmini0832-sourcecharge, def:physics:phyx-mini-0832:phyxminiproblems-problemphyxmini0832-rectangularpointchargesetup}
  The signed Coulomb-law vector produced by one source at a proposed position. Coordinates are metre readouts, charge is read in coulombs, and Physlib's Coulomb constant is read in the corresponding coherent SI units
\end{definition}
\begin{proof}
  This is the declared model definition of coulomb Field Vector At.
\end{proof}
\begin{definition}[Matches Primary Figure832]
  \label{def:physics:phyx-mini-0832:phyxminiproblems-problemphyxmini0832-matchesprimaryfigure832}
  \lean{PhyXMiniProblems.ProblemPhyXMini0832.MatchesPrimaryFigure832}
  \uses{def:physics:phyx-mini-0832:phyxminiproblems-problemphyxmini0832-lengthreadout, def:physics:phyx-mini-0832:phyxminiproblems-problemphyxmini0832-signedchargeincoulombs, def:physics:phyx-mini-0832:phyxminiproblems-problemphyxmini0832-sourcecorner, def:physics:phyx-mini-0832:phyxminiproblems-problemphyxmini0832-rectangularpointchargesetup, def:physics:phyx-mini-0832:phyxminiproblems-problemphyxmini0832-xcoordinateinmeters, def:physics:phyx-mini-0832:phyxminiproblems-problemphyxmini0832-ycoordinateinmeters}
  Literal charge, distance, corner, and presentation readouts from image 832. The bottom-left corner is chosen as the coordinate origin. No resultant field strength or answer choice occurs in this figure predicate
\end{definition}
\begin{proof}
  This is the declared model definition of Matches Primary Figure832.
\end{proof}
\begin{definition}[Uses Standard Coulomb Constant]
  \label{def:physics:phyx-mini-0832:phyxminiproblems-problemphyxmini0832-usesstandardcoulombconstant}
  \lean{PhyXMiniProblems.ProblemPhyXMini0832.UsesStandardCoulombConstant}
  \uses{def:physics:phyx-mini-0832:phyxminiproblems-problemphyxmini0832-rectangularpointchargesetup}
  Standard electrostatic reference data used for the numerical comparison
\end{definition}
\begin{proof}
  This is the declared model definition of Uses Standard Coulomb Constant.
\end{proof}
\begin{definition}[Has Physical Rectangular Charge Parameters]
  \label{def:physics:phyx-mini-0832:phyxminiproblems-problemphyxmini0832-hasphysicalrectangularchargeparameters}
  \lean{PhyXMiniProblems.ProblemPhyXMini0832.HasPhysicalRectangularChargeParameters}
  \uses{def:physics:phyx-mini-0832:phyxminiproblems-problemphyxmini0832-lengthreadout, def:physics:phyx-mini-0832:phyxminiproblems-problemphyxmini0832-rectangularpointchargesetup}
  Positivity and noncoincidence conditions for the physical configuration
\end{definition}
\begin{proof}
  This is the declared model definition of Has Physical Rectangular Charge Parameters.
\end{proof}
\begin{definition}[Satisfies Coulomb Field And Superposition Laws]
  \label{def:physics:phyx-mini-0832:phyxminiproblems-problemphyxmini0832-satisfiescoulombfieldandsuperpositionlaws}
  \lean{PhyXMiniProblems.ProblemPhyXMini0832.SatisfiesCoulombFieldAndSuperpositionLaws}
  \uses{def:physics:phyx-mini-0832:phyxminiproblems-problemphyxmini0832-electricfieldstrengthinnewtonspercoulomb, def:physics:phyx-mini-0832:phyxminiproblems-problemphyxmini0832-sourcecharge, def:physics:phyx-mini-0832:phyxminiproblems-problemphyxmini0832-rectangularpointchargesetup, def:physics:phyx-mini-0832:phyxminiproblems-problemphyxmini0832-resultantmagnitudeinnewtonspercoulombat, def:physics:phyx-mini-0832:phyxminiproblems-problemphyxmini0832-coulombfieldvectorat}
  Coulomb's point-charge field law and vector superposition. The first law is stated at every nonsource point and every time, not only at the target dot. The last field connects the independent dimensionful strength observable to the norm of the resultant vector field; it does not prescribe a number
\end{definition}
\begin{proof}
  This is the declared model definition of Satisfies Coulomb Field And Superposition Laws.
\end{proof}
\begin{definition}[Answer Choice]
  \label{def:physics:phyx-mini-0832:phyxminiproblems-problemphyxmini0832-answerchoice}
  \lean{PhyXMiniProblems.ProblemPhyXMini0832.AnswerChoice}
  Labels attached to the four displayed electric-field choices
\end{definition}
\begin{proof}
  This is the declared model definition of Answer Choice.
\end{proof}
\begin{definition}[displayed Field Strength In Newtons Per Coulomb]
  \label{def:physics:phyx-mini-0832:phyxminiproblems-problemphyxmini0832-displayedfieldstrengthinnewtonspercoulomb}
  \lean{PhyXMiniProblems.ProblemPhyXMini0832.displayedFieldStrengthInNewtonsPerCoulomb}
  \uses{def:physics:phyx-mini-0832:phyxminiproblems-problemphyxmini0832-answerchoice}
  Electric-field strength in N/C printed beside each answer label
\end{definition}
\begin{proof}
  This is the declared model definition of displayed Field Strength In Newtons Per Coulomb.
\end{proof}
\begin{definition}[Is Closest Displayed Field Strength]
  \label{def:physics:phyx-mini-0832:phyxminiproblems-problemphyxmini0832-isclosestdisplayedfieldstrength}
  \lean{PhyXMiniProblems.ProblemPhyXMini0832.IsClosestDisplayedFieldStrength}
  \uses{def:physics:phyx-mini-0832:phyxminiproblems-problemphyxmini0832-electricfieldstrengthquantity, def:physics:phyx-mini-0832:phyxminiproblems-problemphyxmini0832-electricfieldstrengthinnewtonspercoulomb, def:physics:phyx-mini-0832:phyxminiproblems-problemphyxmini0832-answerchoice, def:physics:phyx-mini-0832:phyxminiproblems-problemphyxmini0832-displayedfieldstrengthinnewtonspercoulomb}
  A displayed choice is closest when its printed magnitude is at least as close to the physical resultant strength as every alternative. This models the resolution of the multiple-choice question without defining the physical field strength from the recorded answer
\end{definition}
\begin{proof}
  This is the declared model definition of Is Closest Displayed Field Strength.
\end{proof}
\begin{lemma}[resultant field magnitude bounds]
  \label{lem:physics:phyx-mini-0832:phyxminiproblems-problemphyxmini0832-resultant-field-magnitude-bounds}
  \lean{PhyXMiniProblems.ProblemPhyXMini0832.resultant_field_magnitude_bounds}
  \uses{def:physics:phyx-mini-0832:phyxminiproblems-problemphyxmini0832-electricfieldstrengthinnewtonspercoulomb, def:physics:phyx-mini-0832:phyxminiproblems-problemphyxmini0832-rectangularpointchargesetup, def:physics:phyx-mini-0832:phyxminiproblems-problemphyxmini0832-matchesprimaryfigure832, def:physics:phyx-mini-0832:phyxminiproblems-problemphyxmini0832-usesstandardcoulombconstant, def:physics:phyx-mini-0832:phyxminiproblems-problemphyxmini0832-hasphysicalrectangularchargeparameters, def:physics:phyx-mini-0832:phyxminiproblems-problemphyxmini0832-satisfiescoulombfieldandsuperpositionlaws}
  The three Coulomb contributions give a resultant magnitude between 1.08 * 10\textasciicircum{}5 N/C and 1.09 * 10\textasciicircum{}5 N/C before rounding
\end{lemma}
\begin{proof}
  The conclusion follows from the governing relations and figure readouts named in the statement of resultant field magnitude bounds.
\end{proof}
% --- Archon declaration topology end ---
% --- Archon physics formalization source end ---
